\documentclass[10pt,a4paper]{amsart}
\usepackage{amsmath,amssymb,amsthm,graphicx}
% Retained minimal preamble for the Stacks Project chapter flat.tex (More on Flatness).
% Verbatim excerpt of the chapter preamble: Xy-pic setup (native preamble lines 14-18) and the native
% theorem declarations with the mathematical macros (native preamble lines 154-217). Copied mechanically;
% nothing added, removed or redefined. The wrapper class and the amsmath/amssymb/amsthm/graphicx base
% packages are supplied by the assembler, so the class/externaldocument part of the native preamble is omitted.
% For commutative diagrams we use Xy-pic
\usepackage[all]{xy}

% We use 2cell for 2-commutative diagrams.
\xyoption{2cell}
% Theorem environments.
%
\theoremstyle{plain}
\newtheorem{theorem}[subsection]{Theorem}
\newtheorem{proposition}[subsection]{Proposition}
\newtheorem{lemma}[subsection]{Lemma}

\theoremstyle{definition}
\newtheorem{definition}[subsection]{Definition}
\newtheorem{example}[subsection]{Example}
\newtheorem{exercise}[subsection]{Exercise}
\newtheorem{situation}[subsection]{Situation}

\theoremstyle{remark}
\newtheorem{remark}[subsection]{Remark}
\newtheorem{remarks}[subsection]{Remarks}

\numberwithin{equation}{subsection}

% Macros
%
\def\lim{\mathop{\mathrm{lim}}\nolimits}
\def\colim{\mathop{\mathrm{colim}}\nolimits}
\def\Spec{\mathop{\mathrm{Spec}}}
\def\Hom{\mathop{\mathrm{Hom}}\nolimits}
\def\Ext{\mathop{\mathrm{Ext}}\nolimits}
\def\SheafHom{\mathop{\mathcal{H}\!\mathit{om}}\nolimits}
\def\SheafExt{\mathop{\mathcal{E}\!\mathit{xt}}\nolimits}
\def\Sch{\mathit{Sch}}
\def\Mor{\mathop{\mathrm{Mor}}\nolimits}
\def\Ob{\mathop{\mathrm{Ob}}\nolimits}
\def\Sh{\mathop{\mathit{Sh}}\nolimits}
\def\NL{\mathop{N\!L}\nolimits}
\def\CH{\mathop{\mathrm{CH}}\nolimits}
\def\proetale{{pro\text{-}\acute{e}tale}}
\def\etale{{\acute{e}tale}}
\def\QCoh{\mathit{QCoh}}
\def\Ker{\mathop{\mathrm{Ker}}}
\def\Im{\mathop{\mathrm{Im}}}
\def\Coker{\mathop{\mathrm{Coker}}}
\def\Coim{\mathop{\mathrm{Coim}}}

% Boxtimes
%
\DeclareMathSymbol{\boxtimes}{\mathbin}{AMSa}{"02}

%
% Macros for moduli stacks/spaces
%
\def\QCohstack{\mathcal{QC}\!\mathit{oh}}
\def\Cohstack{\mathcal{C}\!\mathit{oh}}
\def\Spacesstack{\mathcal{S}\!\mathit{paces}}
\def\Quotfunctor{\mathrm{Quot}}
\def\Hilbfunctor{\mathrm{Hilb}}
\def\Curvesstack{\mathcal{C}\!\mathit{urves}}
\def\Polarizedstack{\mathcal{P}\!\mathit{olarized}}
\def\Complexesstack{\mathcal{C}\!\mathit{omplexes}}
% \Pic is the operator that assigns to X its picard group, usage \Pic(X)
% \Picardstack_{X/B} denotes the Picard stack of X over B
% \Picardfunctor_{X/B} denotes the Picard functor of X over B
\def\Pic{\mathop{\mathrm{Pic}}\nolimits}
\def\Picardstack{\mathcal{P}\!\mathit{ic}}
\def\Picardfunctor{\mathrm{Pic}}
\def\Deformationcategory{\mathcal{D}\!\mathit{ef}}

% --- appended 2026-10-12 (Chief): the project own definitions of the reference/slogan/history
% display environments, copied mechanically from the shared Stacks preamble lines 8-12
% (\usepackage{verbatim} plus the three \newenvironment lines, which the chapter preamble keeps
% verbatim from the project preamble). Nothing added, removed or redefined: these environments
% carry no mathematical content and the project itself typesets them as comments. Without these
% three lines a counted statement whose own text uses them fails with Environment slogan undefined.
\usepackage{verbatim}
\newenvironment{reference}{\comment}{\endcomment}
\newenvironment{slogan}{\comment}{\endcomment}
\newenvironment{history}{\comment}{\endcomment}
\usepackage{cleveref}
\crefname{originalproposition}{original proposition}{original propositions}
\Crefname{originalproposition}{Original proposition}{Original propositions}
\makeatletter
\newlabel{algebra-proposition-etale-locally-standard}{{00UE}{source}{}{sourcepointer-1}{}}
\newlabel{algebra-proposition-etale-locally-standard@cref}{{[originalproposition][0][]00UE}{[1][0][]source}{}{}{}}
\makeatother
\crefname{originallemma}{original lemma}{original lemmas}
\Crefname{originallemma}{Original lemma}{Original lemmas}
\makeatletter
\newlabel{lemma-trick-fitting-ideal}{{0812}{source}{}{sourcepointer-2}{}}
\newlabel{lemma-trick-fitting-ideal@cref}{{[originallemma][0][]0812}{[1][0][]source}{}{}{}}
\makeatother
\title{Pushing a finite type ideal along an etale ring map into a product involving an ideal of the base}
\author{The Stacks Project Authors}
\date{stacks-project:flat.tex:local-snapshot-20261011; GNU FDL 1.2 or later}
\begin{document}
\maketitle
\noindent\textit{Editorial note.} All delivered mathematical text is the author's own, copied from the pinned Stacks Project chapter file flat.tex; no statement or proof step was rewritten, reordered or removed.\par\medskip
\noindent\textit{Editorial note.} The preamble is the chapter's own retained minimal preamble (sources/stacks-flat/retained-preamble.tex), a verbatim excerpt of the native Xy-pic setup and of the native theorem declarations and mathematical macros, followed by the project own definitions of the reference, slogan and history display environments (verbatim package plus three newenvironment lines) that the retained file carries at its end and that the project itself typesets as comments. No macro was invented or redefined, and the wrapper class and the amsmath/amssymb/amsthm/graphicx base packages are supplied by the assembler.\par\medskip
\noindent\textit{Editorial note.} The counted statement and proof contain no cross reference, no citation and no included file, so the delivered body is self contained.\par\medskip
\noindent\textit{Editorial note.} Canonical location of the counted proposition: Stacks tag 0813 (sources/stacks-flat/pinned-original-source.tex line 8876).\par\medskip
\noindent\textit{Allowed standard prerequisites (editorial).}\begin{itemize}
\item An etale ring map is standard etale after localising at a suitable element. Algebra, Proposition labelled algebra-proposition-etale-locally-standard (Stacks tag 00UE), invoked at proof lines 8888-8889. Supplies the reduction to B = C\_h with C = A[x]/(f) for a monic f of degree d, which is what makes the exponent dm computable.
\item The fitting ideal trick relating an ideal to a fittable presentation. Same chapter, Lemma labelled lemma-trick-fitting-ideal (source line 8849), invoked at proof line 8896. Supplies the recipe I = Fit\_0(C/(h\_1, ..., h\_\{r+1\})) used to produce the ideal of the base; the verification that it has the required vanishing locus is inside the counted span.
\end{itemize}\medskip
\ifdefined\hypertarget\hypertarget{sourcepointer-1}{}\fi
\noindent\textit{Source reference.} Original proposition 00UE: sources/stacks-chow/tags.txt, source lines 1108--1108. Cross-chapter standard prerequisite located in the project tag list staged in the local snapshot; invoked at proof line 8889 of the counted proof. It is retained only as a navigation pointer; every load-bearing step is inside the delivered span.\par\medskip
\ifdefined\hypertarget\hypertarget{sourcepointer-2}{}\fi
\noindent\textit{Source reference.} Original lemma 0812: sources/stacks-flat/pinned-original-source.tex, source lines 8849--8849. Same-chapter pointer into the author environment lemma at source line 8849; invoked at proof line 8896 of the counted proof. It is retained only as a navigation pointer; every load-bearing step is inside the delivered span.\par\medskip
\begin{lemma}
\label{lemma-push-ideal}
Let $A \to B$ be an \'etale ring map. Let $a \in A$ be a nonzerodivisor.
Let $J \subset B$ be a finite type ideal with $V(J) \subset V(aB)$.
For every $\mathfrak q \subset B$ there exists a finite type ideal
$I \subset A$ with $V(I) \subset V(a)$ and
$g \in B$, $g \not \in \mathfrak q$ such that
$IB_g = JJ'$ for some finite type ideal $J' \subset B_g$.
\end{lemma}

\begin{proof}
We may replace $B$ by a principal localization at
an element $g \in B$, $g \not \in \mathfrak q$. Thus we may assume
that $B$ is standard \'etale, see
Algebra, Proposition \ref{algebra-proposition-etale-locally-standard}.
Thus we may assume $B$ is a localization
of $C = A[x]/(f)$ for some monic $f \in A[x]$ of some degree $d$.
Say $B = C_h$ for some $h \in C$. Choose elements $h_1, \ldots, h_n \in C$
which generate $J$
over $B$. The condition $V(J) \subset V(aB)$ signifies that
$a^m = \sum b_i h_i$ in $B$ for some large $m$. Set $h_{n + 1} = a^m$.
As in Lemma \ref{lemma-trick-fitting-ideal} we take
$I = \text{Fit}_0(C/(h_1, \ldots, h_{r + 1}))$.
Since the module $C/(h_1, \ldots, h_{r + 1})$
is annihilated by $a^m$ we see that $a^{dm} \in I$ which implies
that $V(I) \subset V(a)$.
\end{proof}

\end{document}
